\documentclass{amsart}
% For dealing with references we use the comment environment
\usepackage{verbatim}
\newenvironment{reference}{\comment}{\endcomment}
%\newenvironment{reference}{}{}
\newenvironment{slogan}{\comment}{\endcomment}
\newenvironment{history}{\comment}{\endcomment}

% For commutative diagrams we use Xy-pic
\usepackage[all]{xy}

% We use 2cell for 2-commutative diagrams.
\xyoption{2cell}
\UseAllTwocells

% We use multicol for the list of chapters between chapters
\usepackage{multicol}

% This is generally recommended for better output
\usepackage{lmodern}
\usepackage[T1]{fontenc}

% For cross-file-references

% Package for hypertext links:
\usepackage{hyperref}

% For any local file, say "hello.tex" you want to link to please

% Theorem environments.
%
\theoremstyle{plain}
\newtheorem{theorem}[subsection]{Theorem}
\newtheorem{proposition}[subsection]{Proposition}
\newtheorem{lemma}[subsection]{Lemma}

\theoremstyle{definition}
\newtheorem{definition}[subsection]{Definition}
\newtheorem{example}[subsection]{Example}
\newtheorem{exercise}[subsection]{Exercise}
\newtheorem{situation}[subsection]{Situation}

\theoremstyle{remark}
\newtheorem{remark}[subsection]{Remark}
\newtheorem{remarks}[subsection]{Remarks}

\numberwithin{equation}{subsection}

% Macros
%
\def\lim{\mathop{\mathrm{lim}}\nolimits}
\def\colim{\mathop{\mathrm{colim}}\nolimits}
\def\Spec{\mathop{\mathrm{Spec}}}
\def\Hom{\mathop{\mathrm{Hom}}\nolimits}
\def\Ext{\mathop{\mathrm{Ext}}\nolimits}
\def\SheafHom{\mathop{\mathcal{H}\!\mathit{om}}\nolimits}
\def\SheafExt{\mathop{\mathcal{E}\!\mathit{xt}}\nolimits}
\def\Sch{\mathit{Sch}}
\def\Mor{\mathop{\mathrm{Mor}}\nolimits}
\def\Ob{\mathop{\mathrm{Ob}}\nolimits}
\def\Sh{\mathop{\mathit{Sh}}\nolimits}
\def\NL{\mathop{N\!L}\nolimits}
\def\CH{\mathop{\mathrm{CH}}\nolimits}
\def\proetale{{pro\text{-}\acute{e}tale}}
\def\etale{{\acute{e}tale}}
\def\QCoh{\mathit{QCoh}}
\def\Ker{\mathop{\mathrm{Ker}}}
\def\Im{\mathop{\mathrm{Im}}}
\def\Coker{\mathop{\mathrm{Coker}}}
\def\Coim{\mathop{\mathrm{Coim}}}

% Boxtimes
%
\DeclareMathSymbol{\boxtimes}{\mathbin}{AMSa}{"02}

%
% Macros for moduli stacks/spaces
%
\def\QCohstack{\mathcal{QC}\!\mathit{oh}}
\def\Cohstack{\mathcal{C}\!\mathit{oh}}
\def\Spacesstack{\mathcal{S}\!\mathit{paces}}
\def\Quotfunctor{\mathrm{Quot}}
\def\Hilbfunctor{\mathrm{Hilb}}
\def\Curvesstack{\mathcal{C}\!\mathit{urves}}
\def\Polarizedstack{\mathcal{P}\!\mathit{olarized}}
\def\Complexesstack{\mathcal{C}\!\mathit{omplexes}}
% \Pic is the operator that assigns to X its picard group, usage \Pic(X)
% \Picardstack_{X/B} denotes the Picard stack of X over B
% \Picardfunctor_{X/B} denotes the Picard functor of X over B
\def\Pic{\mathop{\mathrm{Pic}}\nolimits}
\def\Picardstack{\mathcal{P}\!\mathit{ic}}
\def\Picardfunctor{\mathrm{Pic}}
\def\Deformationcategory{\mathcal{D}\!\mathit{ef}}

\setlength{\emergencystretch}{3em}
\begin{document}
\title[Stacks Project, Tag 0C3N]{580 --- Normalization commutes with separable field extension and changes only by a universal homeomorphism under arbitrary field extension}
\maketitle
\noindent Copyright (C) 2005--2025 Johan de Jong. Stacks Project authors. Source: \href{https://stacks.math.columbia.edu/tag/0C3N}{Tag 0C3N}.

\noindent Editorial difficulty note (not author text). Difficulty H2: The proof compares total rings of fractions and their integral closures while allowing nilpotents after inseparable base change. Every integral element is descended to a finitely generated subfield and a suitable Frobenius power is transferred to its maximal separable subextension, proving the universal-homeomorphism assertion..

\noindent Editorial provenance note (not author text). History: excerpt from revision \texttt{a0444\allowbreak 6e57e\allowbreak c1fbc\allowbreak 252a8\allowbreak 71afc\allowbreak ec775\allowbreak 2fb28\allowbreak 07b14}, varieties.tex, lines 5026--5099. Reference presentation was adapted; original-author context and core support blocks were retained or added. The mathematical wording is unchanged. Licensed under GNU FDL 1.2 or later; no invariant sections or cover texts. See ../licenses/COPYING. Context blocks reproduce author definitions and supporting results; they are not additional counted problems.

\begin{lemma}
\label{lemma-normalization-and-change-of-fields}
Let $k$ be a field. Let $X$ be a locally algebraic $k$-scheme.
Let $K/k$ be an extension of fields. Let $\nu : X^\nu \to X$
be the normalization of $X$ and let $Y^\nu \to X_K$ be the
normalization of the base change. Then the canonical morphism
$$
Y^\nu \longrightarrow X^\nu \times_{\Spec(k)} \Spec(K)
$$
is an isomorphism if $K/k$ is separable and a universal homeomorphism
in general.
\end{lemma}

\begin{proof}
Set $Y = X_K$. Let $X^{(0)}$, resp.\ $Y^{(0)}$ be the set of generic points
of irreducible components of $X$, resp.\ $Y$. Then the projection morphism
$\pi : Y \to X$ satisfies $\pi(Y^{(0)}) = X^{(0)}$. This is true because
$\pi$ is surjective, open, and generizing, see
Morphisms, Lemmas \href{https://stacks.math.columbia.edu/tag/0383}{lemma scheme over field universally open (Tag 0383)} and
\href{https://stacks.math.columbia.edu/tag/040F}{lemma open generizing (Tag 040F)}.
If we view $X^{(0)}$, resp.\ $Y^{(0)}$ as (reduced) schemes, then
$X^\nu$, resp.\ $Y^\nu$ is the normalization of $X$, resp.\ $Y$ in
$X^{(0)}$, resp.\ $Y^{(0})$.
Thus Morphisms, Lemma \href{https://stacks.math.columbia.edu/tag/035J}{lemma functoriality normalization (Tag 035J)}
gives a canonical morphism $Y^\nu \to X^\nu$ over $Y \to X$ which in
turn gives the canonical morphism of the lemma by the universal
property of the fibre product.

\medskip\noindent
To prove this morphism has the properties stated in the lemma we may
assume $X = \Spec(A)$ is affine. Let $Q(A_{red})$ be the
total ring of fractions of $A_{red}$. Then $X^\nu$ is the spectrum of
the integral closure $A'$ of $A$ in $Q(A_{red})$, see
Morphisms, Lemmas \href{https://stacks.math.columbia.edu/tag/035O}{lemma normalization reduced (Tag 035O)} and
\href{https://stacks.math.columbia.edu/tag/035P}{lemma description normalization (Tag 035P)}.
Similarly, $Y^\nu$ is the spectrum of the integral closure $B'$ of
$A \otimes_k K$ in $Q((A \otimes_k K)_{red})$. There is a canonical
map $Q(A_{red}) \to Q((A \otimes_k K)_{red})$, a canonical map
$A' \to B'$, and the morphism of the lemma corresponds to the induced map
$$
A' \otimes_k K \longrightarrow B'
$$
of $K$-algebras. The kernel consists of nilpotent
elements as the kernel of $Q(A_{red}) \otimes_k K \to Q((A \otimes_k K)_{red})$
is the set of nilpotent elements.

\medskip\noindent
If $K/k$ is separable, then $A' \otimes_k K$ is normal by
Lemma \href{https://stacks.math.columbia.edu/tag/0C3M}{lemma base change normal by separable (Tag 0C3M)}. In particular
it is reduced, whence $Q((A \otimes_k K)_{red}) = Q(A' \otimes_k K)$
and $B' = A' \otimes_k K$ by
Algebra, Lemma \href{https://stacks.math.columbia.edu/tag/030C}{lemma characterize reduced ring normal (Tag 030C)}.

\medskip\noindent
Assume $K/k$ is not separable. Then the characteristic of $k$ is $p > 0$.
We will show that for every $b \in B'$ there is a power $q$ of $p$
such that $b^q$ is in the image of $A' \otimes_k K$. This will prove
that the displayed map is a universal homeomorphism by
Algebra, Lemma \href{https://stacks.math.columbia.edu/tag/0BRA}{lemma p ring map (Tag 0BRA)}.
For a given $b$ there is a subfield $F \subset K$
with $F/k$ finitely generated such that $b$ is contained in
$Q((A \otimes_k F)_{red})$ and is integral over $A \otimes_k F$.
Choose a monic polynomial $P = T^d + \alpha_1 T^{d - 1} + \ldots + \alpha _d$
with $P(b) = 0$ and $\alpha_i \in A \otimes_k F$.
Choose a transcendence basis $t_1, \ldots, t_r$ for $F$ over $k$.
Let $F/F'/k(t_1, \ldots, t_r)$ be the maximal separable
subextension (Fields,  Lemma \href{https://stacks.math.columbia.edu/tag/030K}{lemma separable first (Tag 030K)}).
Since $F/F'$ is finite purely inseparable, there is a $q$ such
that $\lambda^q \in F'$ for all $\lambda \in F$. Then $b^q$ is
in $Q((A \otimes_k F')_{red})$ and satisfies the polynomial
$T^d + \alpha_1^q T^{d - 1} + \ldots + \alpha _d^q$ with
$\alpha_i^q \in A \otimes_k F'$. By the separable case
we see that $b^q \in A' \otimes_k F'$ and the proof is complete.
\end{proof}
\end{document}
